\chapter{TP \og CLASSIFICATION \fg{} — corrigé complet}

Ce TP (document scanné) demande de \textbf{dessiner deux arbres de décision} et de
\textbf{vérifier avec WEKA et PYTHON}. Les deux exercices reprennent exactement les jeux de
données traités au chapitre~4 ; on en redonne ici la résolution condensée et la procédure de
vérification.

\section{Exercice 1 — Indice de Gini}
\begin{center}\small
\begin{tabular}{cccc}
\toprule
ID & Âge & Type\_voit & Risque\\
\midrule
0 & 23 & Familiale & Élevé\\ 1 & 17 & Sport & Élevé\\ 2 & 43 & Sport & Élevé\\
3 & 68 & Familiale & Faible\\ 4 & 32 & Camion & Faible\\ 5 & 20 & Familiale & Élevé\\
\bottomrule
\end{tabular}
\end{center}

\textbf{Résolution} (détaillée au \S4.2.1) : $GINI(\text{racine})=0{,}444$. Les deux meilleurs
premiers découpages (Type\_voit et Âge $<27{,}5$) sont à égalité ($GINI_{\text{split}}=0{,}222$) ;
scikit-learn retient \textbf{Âge $<27{,}5$}, puis \textbf{Type\_voit} sépare parfaitement.

\begin{Exemple}[arbre attendu]
\texttt{Âge < 27,5 ?} $\rightarrow$ OUI : \textbf{Élevé} ; NON : \texttt{Type\_voit} $\rightarrow$ Sport=\textbf{Élevé},
Familiale=\textbf{Faible}, Camion=\textbf{Faible}. \\
\textbf{Script :} \texttt{uv run python scripts/gini\_arbre.py} (figure \texttt{gini\_arbre.png}).
\end{Exemple}

\section{Exercice 2 — Gain d'information}
Jeu \og achat d'ordinateur \fg{} (14 lignes, \S4.2.2) : $Entropie(S)=0{,}940$. Gains :
$\text{Âge}=0{,}246$ (max), Étudiant $=0{,}151$, État\_civil $=0{,}048$, Revenue $=0{,}029$.

\begin{Exemple}[arbre attendu]
\texttt{Âge ?} $\rightarrow$ $\le30$ : \texttt{Étudiant} (oui=oui, non=non) ; $31..40$ : \textbf{oui} ;
$>40$ : \texttt{État\_civil} (célibataire=oui, marié=non).\\
\textbf{Script :} \texttt{uv run python scripts/gain\_information.py} (figure
\texttt{gain\_information\_arbre.png}).
\end{Exemple}

\section{Vérification sous WEKA}
\begin{enumerate}[leftmargin=1.6em,nosep]
  \item Enregistrer le tableau au format \texttt{.arff} ou \texttt{.csv}.
  \item \emph{Explorer} $\rightarrow$ onglet \emph{Classify}.
  \item Choisir \texttt{trees > J48} (implémentation de C4.5, $\approx$ gain d'information) ou
        \texttt{trees > SimpleCart} (indice de Gini).
  \item \emph{Test options} : \og Use training set \fg{} (jeux minuscules) puis lancer.
  \item Lire l'arbre dans la sortie texte et le comparer à celui de Python — \textbf{ils
        coïncident} (même attribut racine, mêmes feuilles).
\end{enumerate}

\begin{Astuce}
WEKA J48 utilise le \emph{gain ratio} (C4.5) : sur ces petits jeux, le résultat est identique
à l'arbre par gain d'information ci-dessus. Pour Gini, utiliser SimpleCart/CART.
\end{Astuce}

\begin{Resume}
Les deux arbres du TP sont exactement ceux du chapitre~4. Python (scikit-learn) et WEKA
(J48 / SimpleCart) reconstruisent les mêmes structures, ce qui valide la résolution manuelle.
\end{Resume}
